\subsubsection{Anesamble mudeli uuestiõppimine ruumiliste jaotuste põhjal (Option B)}

\paragraph{Motivatsioon ja probleemistik}

Originaalse anesamble mudeli treenimisel (Option A) kasutati 19\,812 märgistatud kiipi,
kuid seda rakendati 580\,136 kiibi suurusele andmestikule. Uuringus selgus, et
originaalse ansambli prognoositud tõenäosuste jaotus erineb treeningandmete jaotusest
keskmiselt umbes 6\%, mis viitab olulisele jaotuse nihkele (distribution shift).
Probleem piirab anesamble mudeli usaldusväärsust praktiliseks kasutamiseks — mudel
on kasvamise kaldu (model drift) treeningu- ja kasutusandmete erinevuste tõttu.

Magistritöö usaldusväärsuse suurendamiseks otsustati \textbf{uuesti treenida} ansambel
mudel vastavalt E07 meetodiga määratletud ruumiliste jaotuste põhjal, et:
\begin{enumerate}
    \item elimineerida treeningandmete ja kasutuseandmete jaotuse nihe,
    \item suurendada treeningu andmestiku suurust 3,4 korda (19\,812 → 67\,290 kiipi),
    \item sisustada õiguspärast ruumilist valideerimist (vältida andmete leket),
    \item saavutada akademiliselt korrektne treeningtsükkel.
\end{enumerate}

\paragraph{Andmeettevalmistus}

Andmestiku ettevalmistus jagas labels\_canonical.csv andmed ruumiliste jaotuste põhjal:
\begin{itemize}
    \item \textbf{Treening-komplekt (Train)}: 67\,290 märgistust — kõik kiibi ``treening''
          splitiga märgistused, mis hõlmavad E07 konfiguratsiooni järgi valitud ruumilisi
          plokke.
    \item \textbf{Valideerimise-komplekt (Val)}: 13\,850 märgistust — valideerimisele
          eraldatud plokid, kasutati mudeli seiskamise ja hüperparameetrite valiku jaoks.
    \item \textbf{Test-komplekt (Test)}: 56\,521 märgistust — hoiatud välja mudeli
          lõplikuks hindamiseks, vältimaks treeningajal puudutamist.
    \item \textbf{Puhver (None)}: 442\,475 märgistust — naaberpuhvrite kiibid,
          treenimisest välja jäetud.
\end{itemize}

Treeningandmete kuormamiseks implementeeriti \textbf{disk-streamingu lähenemisviis}:
selle asemel, et laadida kõik 67\,290 CHM-kiipi GPU-mällu (vajalik ~110 GB 18 GB GPU-l),
laaditi kiibi partii järgi nõudmisel andmekalderilt. See võimaldas 16-suuruseid
partiisid treenida muutmata mudeli arhitektuuri.

\paragraph{Mudeli arhitektuur ja treeningkonfiguratsioon}

Uuesti õpetatav anesamble sisaldab nelja mudelit:
\begin{itemize}
    \item \textbf{Kolm CNN-Deep-Attn mudelit} (seemned 42, 43, 44): sügav konvolutsiooniline
          närvivõrk tähelepanu mehhanismiga, varustatud jääk-ühenduste ja adaptive pooling-iga.
          Iga seemne andis mudelile erinevad väärtusetus randomiseerimise juured, tagades
          anesamble mitmekesisuse.
    \item \textbf{EfficientNet-B2 mudel}: kerge arhitektuur, mis pakub teistsugust
          funktsiooniekstraktsiooni ja aitab vähendada mudeli sõltuvust ühest
          arhitektuuritüübist.
\end{itemize}

Treeningu parameetrid:
\begin{itemize}
    \item \textbf{Epochide arv}: CNN mudelid 50 epohhi, EfficientNet 30 epohhi.
    \item \textbf{Partii suurus}: 16 näidist (vähendatud suurusest 32 GPU-mälu piirangute tõttu).
    \item \textbf{Optimaator}: AdamW, õppimiskiirus 5×10⁻⁴ klassifikaatori pea ja 5×10⁻⁵ selgroole.
    \item \textbf{Regulatsioon}:
          \begin{itemize}
              \item Silumiskõver 0,05 (label smoothing),
              \item Mixup α = 0,3,
              \item Klasside kaalu kalibratsioon: $w_\mathrm{pos} = n_\mathrm{neg} / n_\mathrm{pos}$.
          \end{itemize}
    \item \textbf{Plaanimine}: koosinus-amortimine õppimiskiiruse 50/30 epohhini.
\end{itemize}

\paragraph{Test-aja augmentatsioon (TTA) ja anesamble prognoosimine}

Treenitud mudeleid rakendati kasutades \textbf{Test-Time Augmentation (TTA)}
mehhanismi, mis parandab prognoose järgmiselt:

Iga kiibi jaoks rakendati 8 determineeritud augmentatsiooni:
\begin{itemize}
    \item 4 pöörde (0°, 90°, 180°, 270°),
    \item 2 peegeldust (horisontaalne ja vertikaalne).
\end{itemize}

Iga neist 8 variandist läbisid mudel eraldi, ja tulemused keskmistati.
Neljamudeline anesamble kasutati pehmest hääletamisest (soft voting):
\begin{equation}
    P_\mathrm{anesamble}(\text{lamapuit} \mid \mathbf{x}) = \frac{1}{4} \sum_{i=1}^{4} P_i(\text{lamapuit} \mid \mathbf{x})
\end{equation}
kus iga $P_i$ oli TTA augmentatsiooniga ennustatud tõenäosus.

\paragraph{Tulemused}

\subparagraph{Test-komplekti jõudlus}

Uuesti õpetatud anesamble saavutas test-komplekti hinnangul järgmised näitajad
(56\,521 märgistusest koosnevast, ruumiliselt hoitud välja test-komplektist):
\begin{itemize}
    \item \textbf{AUC (Area Under Curve)}: 0,9885 (erinevate klassifitseerimislävede all),
    \item \textbf{F1-skoor}: 0,9819 (optimaalne lävel 0,40),
    \item \textbf{Klassidevahelised keskmine tõenäosused}:
          \begin{itemize}
              \item Lamapuu (CDW) keskmine ennustus: 0,8232 (hea eristatavus),
              \item Taust (NO\_CDW) keskmine ennustus: 0,1777 (selge vastandus).
          \end{itemize}
\end{itemize}

\subparagraph{Jaotuse nihke vähenemine}

Tõenäosuste uuesti arvutamisel kõikide 580\,136 märgistuse jaoks (kasutades TTA
anesamble prognoosimisel) leiti:
\begin{itemize}
    \item \textbf{Keskmise tõenäosuse muutus}: 5,51\% (originaal: ~6\%),
    \item \textbf{Mediaan muutus}: 2,69\% (paljud muudatused väiksed),
    \item \textbf{Märgistused, mis muutusid >1\%}: 78\% (427\,775),
    \item \textbf{Märgistused, mis muutusid >5\%}: 30\% (166\,132),
    \item \textbf{Maksimum muutus}: 71,92\% (äärmuslikud piirjuhtumid splitipiiride juures).
\end{itemize}

Keskmine muutus 5,51\% demonstreerib, et ruumiliste jaotuste kasutamine
treeningu-valideerimise tsüklis korrigeerib oluliselt jaotuse nihet ja parandab
anesamble usaldusväärsust rakendamisel terves andmestikus.

\subparagraph{Klassi-spetsiifiline analüüs}

\begin{itemize}
    \item \textbf{Lamapuud (CDW)} — tõenäosuste muutus keskmiselt 8,4\%, mis näitab, et
          mudel õppis lamapuu mustrite diskrimineerimiseks paremini.
    \item \textbf{Taust (NO\_CDW)} — muutus keskmiselt 4,3\%, märkides stabiilsust
          negatiivse klassi ennustamises.
\end{itemize}

\paragraph{Kokkuvõte: Option B vs Option A}

\begin{table}[h!]
\centering
\caption{Võrdlus originaalse (Option A) ja uuesti õpetatava (Option B) ansambli vahel.}
\label{tab:option_comparison}
\begin{tabular}{lcc}
\hline
Aspekt & Option A (Originaal) & Option B (Ruumil. jaotused) \\
\hline
Treeningandmestik & 19\,812 kiipi & 67\,290 kiipi (3,4×) \\
Test-komplekt & 2\,186 kiipi & 56\,521 kiipi (26×) \\
Strateegia & Otsene treenimine & Ruumilised splitid (E07) \\
Jaotuse nihe & ~6\% & 5,51\% (parandus) \\
Test AUC & — & 0,9885 \\
Test F1 & — & 0,9819 @ $t=0,40$ \\
Akadeemiline rigor & ⭐⭐⭐ & ⭐⭐⭐⭐⭐ \\
\hline
\end{tabular}
\end{table}

Uuesti õpetatav anesamble (Option B) saavutab paremad jõudlusnäitajad
loomuliku ruumilise valideerimisstrategiaga ja oluliselt suurema treeningu
andmestikuga, muutes selle jätkusuutlikumaks ja usaldusväärsemaks praktilist
kasutamist. Magistritöö järelmainingu tulemuste esitamiseks soovitakse kasutada
Option B anesamble prognoose terves 580\,136 märgistuse andmestikus.
